% Generated draft: requires mathematical and source review.
\documentclass[11pt]{article}
\usepackage[margin=25mm]{geometry}
\usepackage{fontspec}
\setmainfont{Latin Modern Roman}
\usepackage{amsmath,amssymb,mathtools,unicode-math}
\setmathfont{Latin Modern Math}
\usepackage{enumitem,needspace,xurl,hyperref}
\setlength{\parindent}{0pt}
\setlength{\parskip}{5pt}
\setlength{\emergencystretch}{3em}
\newcommand{\N}{\mathbb{N}}
\newcommand{\Z}{\mathbb{Z}}
\newcommand{\Q}{\mathbb{Q}}
\newcommand{\R}{\mathbb{R}}
\newcommand{\C}{\mathbb{C}}
\newcommand{\F}{\mathbb{F}}
\newcommand{\A}{\mathbb{A}}
\newcommand{\PP}{\mathbb{P}}
\newcommand{\E}{\mathbb{E}}
\newcommand{\eps}{\varepsilon}
\newcommand{\dd}{\,\mathrm{d}}
\begin{document}
\section*{UnsolvedMath ID: 800716}
\textbf{Source problem code:} catalogue entry
\subsection{Definitions and mathematical statement}
For a finite-dimensional central simple algebra \(A\) over a field \(F\), its degree is \(\deg(A)=\sqrt{\dim_F A}\). Writing \(A\cong M_r(D)\), where \(D\) is a central division algebra over \(F\), define its index by \(\operatorname{ind}(A)=\deg(D)\).

The supplied record does not specify the field hypotheses or the convention for symbol algebras. Under the usual Kummer convention, let \(n\geq 2\) be an integer, let \(F\) be a field with \(\operatorname{char}(F)\nmid n\), and fix a primitive \(n\)-th root of unity \(\zeta\in F\). For \(u,v\in F^\times\), write
\[
(u,v)_{n,\zeta}=F\langle x,y\rangle/(x^n-u,\ y^n-v,\ yx-\zeta xy).
\]
This is a central simple \(F\)-algebra of degree \(n\). Suppress \(\zeta\) from the notation while keeping it fixed. For this formulation, two degree-\(n\) central simple algebras share a common maximal subfield if there is a field extension \(L/F\) with \([L:F]=n\) and unital \(F\)-algebra embeddings of \(L\) into both algebras.

\textbf{Question under this convention.} For every such \(n,F,\zeta\) and every \(a,b,c,d\in F^\times\), does
\[
\operatorname{ind}\bigl((a,b)_n\otimes_F(c,d)_n\bigr)<n^2
\]
imply that there exist \(a',b',c',d'\in F^\times\), a field extension \(L/F\) of degree \(n\), and unital embeddings
\[
L\hookrightarrow(a',b')_n,
\qquad
L\hookrightarrow(c',d')_n,
\]
such that
\[
(a,b)_n\otimes_F(c,d)_n
\cong_F
(a',b')_n\otimes_F(c',d')_n?
\]
Here \(\cong_F\) means isomorphism of \(F\)-algebras. No assumption that any symbol factor is a division algebra is stated in the supplied record. The field hypotheses and the degree-\(n\) interpretation of ``common maximal subfield'' above are explicit conventions for this draft; their agreement with the intended source formulation requires review. In particular, a common field extension is not being replaced by a common commutative \(\acute{e}\)tale algebra. The source does not specify the range of \(n\); under the same convention, \(n=1\) cannot satisfy the strict index inequality.
\subsection{Short English statement}
The question asks whether a tensor product of two degree-\(n\) symbol algebras whose index is less than \(n^2\) can always be expressed as a tensor product of two degree-\(n\) symbol algebras having a common maximal subfield. The symbol convention, field hypotheses, and meaning of maximal subfield require confirmation. The registry labels the problem open; that label is supplied source data, not an independently verified assessment.
\subsection{Sources}
\begin{itemize}
\item Local UnsolvedMath registry, permanent ID \(800716\), exact source problem code \texttt{catalogue entry}; title: \emph{Open Problems on Central Simple Algebras --- Problem 30}. The supplied record attributes the problem to Auel, Brussel, Garibaldi, and Vishne and gives the proposed year as 2010.
\item Auel, Brussel, Garibaldi, and Vishne, \emph{Open Problems on Central Simple Algebras} (2010), Problem 30. Supplied source URL: \url{https://arxiv.org/abs/1006.3304}.
\item The supplied provenance reports access on 2026-10-16 and extraction by \texttt{arxiv\_latex}. These are supplied provenance statements; the source has not been independently accessed or verified for this draft.
\item The supplied registry record is labeled \texttt{open} and was updated on 2026-10-16. This is dated registry curation, not an independently verified literature assessment. No separate \texttt{literature\_sources} list or \texttt{literature\_checked\_at} date was supplied.
\end{itemize}
\end{document}
